\section{Conclusion}

We presented Dynamic Speculative Planning (DSP), a method that addresses fixed-step speculative execution limitations in LLM-based agents through lightweight online reinforcement learning requiring no pre-deployment preparation. Our evaluation demonstrates DSP reduces redundant computational overhead by up to 60\% and overall cost by 30\% while maintaining acceleration benefits, with user-controllable parameters enabling precise calibration of latency-cost tradeoffs. 
This advancement improves the viability of deploying sophisticated agents in latency-sensitive real-world applications.

%The key innovation of our approach is the introduction of user-controllable acceleration through a single hyperparameter that navigates the latency-cost continuum, enabling practitioners to precisely calibrate system behavior according to their specific operational constraints. This flexibility is increasingly valuable in the rapidly evolving LLM ecosystem with fluctuating pricing structures and computational capabilities. DSP represents a significant step toward making LLM-based agents more practical for real-world deployment by intelligently balancing performance and cost efficiency.
%\textbf{Limitations and Future Work}
%Despite the promising results of Dynamic Speculative Planning, several limitations remain. Our framework does not adaptively select or optimize the approximation agent itself, though the speed, accuracy, and availability characteristics of this component significantly impact overall system performance. This represents an important direction for future work. Additionally, our current verification mechanism employs relatively strict rejection criteria that may incorrectly reject valid alternative approaches when they differ syntactically from the target agent's preferred solution. This high false negative rate unnecessarily constrains parallelization opportunities, suggesting that semantic verification mechanisms recognizing functionally equivalent actions could further enhance acceleration. While our online learning approach eliminates pre-deployment requirements, techniques to accelerate the initial learning phase and more extensive evaluation across diverse domains would further strengthen our approach's practical utility.